\documentclass[10pt,a4paper]{amsart}
\usepackage[margin=19mm]{geometry}
\usepackage{amsmath,amssymb,mathtools}
\usepackage[scr=rsfs,cal=euler]{mathalfa}
\usepackage{xfrac}
\usepackage[unicode,hidelinks,bookmarks=false]{hyperref}
\newcommand{\ie}{i.e.\ }
\newcommand{\cf}{cf.\ }
\newcommand{\resp}{respectively\ }
\newcommand{\Subsection}{\subsection}
\providecommand{\nobreakdash}{\nobreak\hbox{-}}
\setlength{\emergencystretch}{2em}
\multlinegap0pt
\usepackage{float}
\usepackage{amssymb}
\usepackage{bm}
\usepackage{mathtools}
\usepackage{braket}
\usepackage{enumerate}
\usepackage{tikz}
\usepackage{tcolorbox}
\newtheorem{oss}{Remark}
\newtheorem{lemma}{Lemma}
\title{On the Effectless Cut Method for Laplacian Eigenvalues in any dimensions}
\author{Vincenzo Amato, Nunzia Gavitone, Francesca de Giovanni}
\date{Source: arXiv:2604.00976v1 (CC BY 4.0)}
\begin{document}
\maketitle

\noindent\emph{Source and licence.} On the Effectless Cut Method for Laplacian Eigenvalues in any dimensions. Version arXiv:2604.00976v1 (CC BY 4.0), distributed by arXiv under the Creative Commons Attribution 4.0 International licence, http://creativecommons.org/licenses/by/4.0/, as recorded in the arXiv licence metadata for this exact version and shown on the abstract page https://arxiv.org/abs/2604.00976v1. Full licence notice: \texttt{licenses/arxiv-2604.00976-CC-BY-4.0.txt}.\par\smallskip

\noindent\textit{Editorial scope.} Counted unit: the article's lemma G\_aperti (source0.tex lines 485--493). The article's complete original proof of it is delivered with it (source0.tex lines 495--561, one \texttt{proof} environment of 67 lines). No mathematics is written, repaired or re-derived: the statement and the proof are the article's own lines, in the article's own notation, and no retained line is substituted or re-flowed.  Retained uncounted as the source-local context the proof needs: lines 277--284; lines 405--420; lines 473--476; lines 478--481.  Nothing was omitted from any retained span, so the package carries no omission record.  No retained line contains a citation, so the wrapper carries no bibliography and no bibliography entry.  No retained line contains a figure environment, an image-inclusion command or drawing code, so no image is delivered and the package declares no asset dependency.  Editorial formatting only, in addition: an AMS \texttt{amsart} preamble that restates the article's own declarations and macros used by the retained lines, namely the environments \texttt{oss}, \texttt{lemma} on separate counters, together with the standard packages those lines need.  The retained environments print in this document as Oss 1, Lemma 1 in order of appearance, whereas the article prints the same items with the numbering of its own full document.  The complete native source of the article as distributed in the arXiv e-print for this exact version is bundled unchanged as \texttt{sources/197625/source0.tex}.
\begin{equation}
    \label{probprinc}
    \begin{cases}
    -\Delta u = \lambda^{RR}(\Omega)u \qquad & \text{in } \Omega, \\
     \displaystyle \frac{\partial u }{\partial \nu} + \beta_{in} u = 0  \qquad & \text{on } \partial\Omega_{in},\\
    \displaystyle \frac{\partial u }{\partial \nu} + \beta_{out} u = 0  \qquad & \text{on } \partial\Omega_{out}.
\end{cases}
\end{equation}

\begin{oss}\label{hopf}
Thanks to the boundary conditions in~\eqref{probprinc} and the regularity of $u$,
Hopf's lemma yields the following sign properties for the normal derivative:
\begin{itemize}
    \item if $(\beta_{\mathrm{in}}, \beta_{\mathrm{out}}) \in (0,+\infty]^2$, then
    \[
    \frac{\partial u}{\partial \nu}(x) < 0
    \qquad \text{on } \partial \Omega;
    \]
    \item if $(\beta_{\mathrm{in}}, \beta_{\mathrm{out}}) \in (-\infty,0)^2$, then
    \[
    \frac{\partial u}{\partial \nu}(x) > 0
    \qquad \text{on } \partial \Omega.
    \]
\end{itemize}
\end{oss}

\begin{equation}
\label{gin}
    G_{in} = \{ x \in \Omega \mid \exists\, t_x \text{ such that } z_x(t_x) \in \partial \Omega_{in} \},
\end{equation}

\begin{equation}
\label{gout}
    G_{out} = \{ x \in \Omega \mid \exists\, t_x \text{ such that } z_x(t_x) \in \partial \Omega_{out} \}.
\end{equation}

\begin{lemma}\label{G_aperti}
Let $G_{in}$ and $G_{out}$ be defined as in \eqref{gin} and \eqref{gout}. Then they are open, disjoint, and connected sets.
Moreover,
\[
\partial \Omega_{in} \subset G_{in}
\qquad \text{and} \qquad
\partial \Omega_{out} \subset G_{out}.
\]
\end{lemma}

\begin{proof}
We prove the statement for $G_{in}$; the proof for $G_{out}$ follows by analogous arguments.

Since $u \in C^{1}(\overline \Omega)$ and $\partial u / \partial \nu \neq 0$ on $\partial \Omega_{in}$, there exists $\overline{\delta} > 0$ such that
\[
(\partial \Omega_{in})_{\overline{\delta}} \subset G_{in}.
\]

Let $x \in G_{in}$. Then there exists $t_x$ such that
\[
z_x(t_x) = \bar{x} \in \partial \Omega_{in}.
\]
Define
\[
\gamma := \{ z_x(t) \mid t \in [0,t_x] \},
\qquad
m := \inf_{z \in \gamma} |\nabla u(z)| > 0.
\]

Let $y \in B_x(a\delta)$ for some $0 < a < 1$, to be fixed later, and choose
\[
\delta < \min \{ m, \overline{\delta}/8 \}.
\]
Our goal is to show that $y \in G_{in}$, namely that there exists $\tilde{t}$ such that
\[
z_y(\tilde {t}) \in (\partial \Omega_{in})_{\overline{\delta}}.
\]

Let $\bar{t}$ be such that
\[
z_x(\bar{t}) \in (\partial \Omega_{in})_{\overline{\delta}/2}
\setminus (\partial \Omega_{in})_{\overline{\delta}/4}.
\]
Since $u$ is analytic on compact subsets of $\Omega$, there exists $L > 0$ such that
\[
|\nabla u(z) - \nabla u(y)| \le L |z-y|
\qquad \forall z,y \in \Omega_{\overline{\delta}/8}= \left\{x \in \Omega : \, d (x, \partial \Omega )> \frac{\overline \delta
}{8}\right\}.
\]

We compute
\begin{align*}
|z_x(t) - z_y(t)|
&= \left| x - y - \int_0^t \bigl( \nabla u(z_x(s)) - \nabla u(z_y(s)) \bigr)\, ds \right| \\
&\le |x-y| + \int_0^t |\nabla u(z_x(s)) - \nabla u(z_y(s))|\, ds \\
&\le a\delta + L \int_0^t |z_x(s) - z_y(s)|\, ds.
\end{align*}
By Gronwall's lemma, it follows that
\[
|z_x(t) - z_y(t)| \le a \delta e^{Lt}
\qquad \forall t \in [0,\bar{t}].
\]

Choosing $a = e^{-L\bar{t}}$, we obtain
\[
|z_x(\bar{t}) - z_y(\bar{t})| < \delta,
\]
and therefore
\[
z_y(\bar{t}) \in (\partial \Omega_{in})_{\overline{\delta}},
\]
which implies $y \in G_{in}$.

Hence, we have proved that $G_{in}$ contains $\partial \Omega_{in}$ and is open. Since every point $x \in G_{in}$ can be connected to $\partial \Omega_{in}$ (which is connected), it follows that $G_{in}$ is also connected. 

Moreover, we emphasize that no flow line can connect both $\partial \Omega_{in}$ and $\partial \Omega_{out}$; otherwise, one would obtain a path from $\partial \Omega_{in}$ to $\partial \Omega_{out}$ along which the gradient does not change sign, contradicting Remark~\ref{hopf}. This ensures that $G_{in}$ and $G_{out}$ are disjoint.
\end{proof}

\end{document}
